% archon:physics
% archon:covers ArchonPhysics/MicroscopicDynamics.lean

\chapter{Finite microscopic Hamiltonian dynamics}
\label{ch:ArchonPhysics_MicroscopicDynamics}

This chapter fixes the deterministic finite periodic chain associated with the
Hamiltonian of Chapter~\ref{ch:ArchonPhysics_HamiltonianScaling}.  Its phase
space is ordered as position--momentum pairs $(q,p)$.  The main analytic result
is Picard--Lindelof existence on a nontrivial time interval together with
uniqueness only as a germ at the initial time.  The conservation theorems are
conditional: the concrete ODE must hold at every point of the unoriented
interval between the two endpoints.  In particular, the global-solution object
below is only a predicate; this chapter proves no global existence, random or
thermodynamic limit, kinetic limit, equilibration, or thermalization theorem.

\begin{definition}\leanok
[Position--momentum phase space]
\label{def:microscopic:phase-space}
\lean{ArchonPhysics.MicroscopicDynamics.PhaseSpace}
\uses{def:lattice:configuration}
For a periodic chain with $N$ sites, set
\[
  \mathcal P_N=\operatorname{Configuration}_N\times
  \operatorname{Configuration}_N,
\]
with a state written $z=(q,p)$.
\end{definition}
\begin{proof}

\leanok
This is the product type of two finite real lattice configurations, with the
first projection designated as position and the second as momentum.
\end{proof}

\begin{definition}\leanok
[Single-bond interaction potential]
\label{def:microscopic:interaction-potential}
\lean{ArchonPhysics.MicroscopicDynamics.interactionPotential}
For degree $n\in\mathbb N$ and coupling $\lambda\in\mathbb R$, define
\[
  V_{n,\lambda}(x)=\frac{x^2}{2}+\frac{\lambda}{n}x^n.
\]
The Lean expression is total for every natural $n$; derivative results below
use the guard $1\le n$.
\end{definition}
\begin{proof}

\leanok
This is the polynomial bond energy occurring termwise in the locked finite
Hamiltonian.
\end{proof}

\begin{definition}\leanok
[Single-bond interaction force]
\label{def:microscopic:interaction-force}
\lean{ArchonPhysics.MicroscopicDynamics.interactionForce}
Define the polynomial expression
\[
  F_{n,\lambda}(x)=x+\lambda x^{n-1},
\]
where subtraction in the exponent is natural-number subtraction.  For
$1\le n$, this is the derivative of $V_{n,\lambda}$.
\end{definition}
\begin{proof}

\leanok
This is a definition for all natural degrees.  The degree guard needed to
identify it with the derivative is imposed by the energy-conservation theorem.
\end{proof}

\begin{definition}\leanok
[Periodic nearest-neighbour lattice force]
\label{def:microscopic:lattice-force}
\lean{ArchonPhysics.MicroscopicDynamics.latticeForce}
\uses{def:microscopic:interaction-force,
      def:lattice:forward-difference}
For a configuration $q$, define the force at site $i$ by applying the periodic
forward difference to the shifted bond-force field.  Equivalently, it is the
outgoing bond force minus the incoming bond force.
\end{definition}
\begin{proof}

\leanok
The definition is
$D\bigl(i\mapsto F_{n,\lambda}((Dq)_{i-1})\bigr)$ on the cyclic site set.
\end{proof}

\begin{definition}\leanok
[Microscopic Hamiltonian vector field]
\label{def:microscopic:vector-field}
\lean{ArchonPhysics.MicroscopicDynamics.microscopicVectorField}
\uses{def:microscopic:phase-space, def:microscopic:lattice-force,
      def:lattice:mass-actions}
For a positive mass profile $m$, set
\[
  X_{m,n,\lambda}(q,p)
   =\bigl(M^{-1}p,\,\operatorname{latticeForce}_{n,\lambda}(q)\bigr).
\]
Thus the two components encode $\dot q=M^{-1}p$ and the finite periodic force
law for $\dot p$.
\end{definition}
\begin{proof}

\leanok
This is the ordered pair of the inverse-mass action and the lattice force.
\end{proof}

\begin{definition}\leanok
[Classical solution at a time]
\label{def:microscopic:classical-solution-at}
\lean{ArchonPhysics.MicroscopicDynamics.IsClassicalSolutionAt}
\uses{def:microscopic:vector-field}
A path $z:\mathbb R\to\mathcal P_N$ is a classical solution at time $t$ when
\[
  z'(t)=X_{m,n,\lambda}(z(t))
\]
in the Fr\'echet-derivative sense encoded by \texttt{HasDerivAt}.
\end{definition}
\begin{proof}

\leanok
This predicate is exactly the displayed vector-valued derivative equation at
the specified time.
\end{proof}

\begin{definition}\leanok
[Global classical-solution predicate]
\label{def:microscopic:global-classical-solution}
\lean{ArchonPhysics.MicroscopicDynamics.IsGlobalClassicalSolution}
\uses{def:microscopic:classical-solution-at}
A path is called a global classical solution when the classical ODE predicate
holds at every real time.
\end{definition}
\begin{proof}

\leanok
This is universal quantification of the pointwise solution predicate over
$\mathbb R$.  It does not assert that any such path exists.
\end{proof}

\begin{definition}\leanok
[Hamiltonian energy on phase space]
\label{def:microscopic:hamiltonian-energy}
\lean{ArchonPhysics.MicroscopicDynamics.hamiltonianEnergy}
\uses{def:microscopic:phase-space,
      def:hamiltonian:lattice-hamiltonian}
For $z=(q,p)$, define
\[
  E_{m,n,\lambda}(z)=H_{m,n,\lambda}(p,q).
\]
The argument reversal records that phase space is ordered $(q,p)$ whereas the
previous Hamiltonian API is ordered $(p,q)$.
\end{definition}
\begin{proof}

\leanok
Evaluate the existing finite lattice Hamiltonian on the momentum and position
projections in that order.
\end{proof}

\begin{definition}\leanok
[Total canonical momentum]
\label{def:microscopic:total-momentum}
\lean{ArchonPhysics.MicroscopicDynamics.totalMomentum}
\uses{def:microscopic:phase-space}
For $z=(q,p)$, define the total canonical momentum by
\[
  P(z)=\sum_{i\in\mathbb Z/N\mathbb Z}p_i.
\]
\end{definition}
\begin{proof}

\leanok
This is the finite sum of the momentum component over all cyclic sites.
\end{proof}

\begin{theorem}\leanok
[Pointwise lattice-force expansion]
\label{thm:microscopic:lattice-force-apply}
\lean{ArchonPhysics.MicroscopicDynamics.latticeForce_apply}
\uses{def:microscopic:lattice-force,
      def:microscopic:interaction-force}
For every site $i$,
\[
 \operatorname{latticeForce}_{n,\lambda}(q)_i
 =F_{n,\lambda}((Dq)_i)-F_{n,\lambda}((Dq)_{i-1}).
\]
\end{theorem}
\begin{proof}

\leanok
Unfold the two periodic forward differences in the lattice-force definition.
\end{proof}

\begin{theorem}\leanok
[Cancellation of internal forces]
\label{thm:microscopic:sum-lattice-force}
\lean{ArchonPhysics.MicroscopicDynamics.sum_latticeForce}
\uses{def:microscopic:lattice-force,
      thm:lattice:sum-forward-difference}
For a nonempty finite periodic chain,
\[
  \sum_i \operatorname{latticeForce}_{n,\lambda}(q)_i=0.
\]
\end{theorem}
\begin{proof}

\leanok
The force is a periodic forward difference, whose finite sum telescopes to
zero on the cyclic site set.
\end{proof}

\begin{theorem}\leanok
[Continuous differentiability of the vector field]
\label{thm:microscopic:vector-field-contdiff}
\lean{ArchonPhysics.MicroscopicDynamics.microscopicVectorField_contDiff}
\uses{def:microscopic:vector-field,
      def:microscopic:interaction-force,
      def:microscopic:lattice-force}
For a nonempty finite chain, $X_{m,n,\lambda}$ is continuously differentiable
of order one for every natural degree $n$ and every real coupling $\lambda$.
\end{theorem}
\begin{proof}

\leanok
After expanding inverse-mass action and periodic differences, every coordinate
of the finite-dimensional vector field is a polynomial expression composed
with linear coordinate maps.  Such a map is $C^1$.
\end{proof}

\begin{theorem}\leanok
[Existence and uniqueness of a local solution germ]
\label{thm:microscopic:local-solution-germ}
\lean{ArchonPhysics.MicroscopicDynamics.exists_unique_local_solution_germ}
\uses{def:microscopic:phase-space, def:microscopic:vector-field,
      thm:microscopic:vector-field-contdiff}
For every initial time $t_0$ and state $z_0$, there are $\varepsilon>0$ and a
path $z:\mathbb R\to\mathcal P_N$ with $z(t_0)=z_0$ that solves the ODE in the
within-derivative sense on
$[t_0-\varepsilon,t_0+\varepsilon]$.  If another path $y$ has the same initial
value and satisfies the ordinary derivative equation throughout the open
interval, then $z$ and $y$ agree eventually in a neighbourhood of $t_0$.
\end{theorem}
\begin{proof}

\leanok
The $C^1$ vector field is locally Lipschitz, so finite-dimensional
Picard--Lindelof gives a solution on a nontrivial closed interval.  Continuity
keeps both candidate paths in a common local closed ball near $t_0$; local ODE
uniqueness then identifies their germs.  The conclusion is deliberately local
and does not extend the solution to all real times.
\end{proof}

\begin{theorem}\leanok
[Pointwise conservation law for total momentum]
\label{thm:microscopic:momentum-derivative-zero}
\lean{ArchonPhysics.MicroscopicDynamics.totalMomentum_hasDerivAt_zero_at}
\uses{def:microscopic:classical-solution-at,
      def:microscopic:total-momentum,
      thm:microscopic:sum-lattice-force}
Whenever $z$ satisfies the microscopic ODE at $t$, the scalar path
$s\mapsto P(z(s))$ has derivative zero at $t$.
\end{theorem}
\begin{proof}

\leanok
Differentiate the finite sum of momentum coordinates.  The ODE replaces each
coordinate derivative by the lattice force, and the internal-force sum
cancels by periodicity.
\end{proof}

\begin{theorem}\leanok
[Pointwise conservation law for Hamiltonian energy]
\label{thm:microscopic:energy-derivative-zero}
\lean{ArchonPhysics.MicroscopicDynamics.hamiltonianEnergy_hasDerivAt_zero_at}
\uses{def:microscopic:interaction-potential,
      def:microscopic:interaction-force,
      def:microscopic:classical-solution-at,
      def:microscopic:hamiltonian-energy,
      thm:microscopic:lattice-force-apply}
Assume $1\le n$.  Whenever $z$ satisfies the microscopic ODE at $t$, the
scalar path $s\mapsto E_{m,n,\lambda}(z(s))$ has derivative zero at $t$.
\end{theorem}
\begin{proof}

\leanok
Differentiate kinetic and bond-potential energies site by site.  Positivity of
each mass justifies the inverse-mass algebra, and $1\le n$ identifies the
potential derivative with $F_{n,\lambda}$.  The kinetic power and bond power
cancel after a cyclic index shift, leaving derivative zero.
\end{proof}

\begin{theorem}\leanok
[Conditional energy conservation between endpoints]
\label{thm:microscopic:energy-endpoints}
\lean{ArchonPhysics.MicroscopicDynamics.hamiltonianEnergy_eq_of_forall_mem_uIcc}
\uses{def:microscopic:classical-solution-at,
      def:microscopic:hamiltonian-energy,
      thm:microscopic:energy-derivative-zero}
Assume $1\le n$.  If $z$ satisfies the microscopic ODE at every
$u\in\operatorname{uIcc}(s,t)$, then
\[
  E_{m,n,\lambda}(z(s))=E_{m,n,\lambda}(z(t)).
\]
\end{theorem}
\begin{proof}

\leanok
The preceding theorem gives zero within-derivative throughout the unoriented
closed interval.  The mean-value bound on this convex interval makes the
endpoint difference have norm at most zero.  This is conditional conservation;
it does not supply a solution spanning arbitrary prescribed endpoints.
\end{proof}

\begin{theorem}\leanok
[Conditional momentum conservation between endpoints]
\label{thm:microscopic:momentum-endpoints}
\lean{ArchonPhysics.MicroscopicDynamics.totalMomentum_eq_of_forall_mem_uIcc}
\uses{def:microscopic:classical-solution-at,
      def:microscopic:total-momentum,
      thm:microscopic:momentum-derivative-zero}
If $z$ satisfies the microscopic ODE at every
$u\in\operatorname{uIcc}(s,t)$, then
\[
  P(z(s))=P(z(t)).
\]
\end{theorem}
\begin{proof}

\leanok
Total momentum has zero within-derivative throughout the convex unoriented
interval.  The corresponding mean-value bound forces equality at its two
endpoints.  As above, the theorem assumes rather than constructs a path on the
whole interval.
\end{proof}
